% TOP_PROBLEM: {"id": 30006984, "title": "Bodirsky-Pinsker Tractability Conjecture for Infinite-Domain Constraint Satisfaction", "category_id": 15, "category": {"id": 15, "name": "computer_science", "display_name": "Computer Science", "description": "Computational complexity, algorithms, and theoretical CS.", "slug": "computer-science", "order_index": 15, "created_at": "2026-07-31T15:26:25.671Z"}, "status": "open"}
% FIRST_POSTED: 2026-09-25
% ENTRY_KIND: statement_only
% !TeX program = lualatex
% Definition and sources only; this file is not an LLM solution attempt.
\documentclass[11pt]{article}
\usepackage[margin=25mm]{geometry}
\usepackage{fontspec}
\setmainfont{Latin Modern Roman}
\usepackage{amsmath,amssymb,mathtools}
\usepackage{unicode-math}
\setmathfont{Latin Modern Math}
\usepackage{xurl,hyperref}
\hypersetup{colorlinks=true,urlcolor=blue}
\setcounter{secnumdepth}{0}
\setlength{\parindent}{0pt}
\setlength{\parskip}{5pt}
\setlength{\emergencystretch}{3em}
\begin{document}
\section*{\texorpdfstring{Bodirsky-Pinsker Tractability Conjecture for Infinite-Domain Constraint Satisfaction}{Bodirsky-Pinsker Tractability Conjecture for Infinite-Domain Constraint Satisfaction}}\label{30006984}
\subsection{Definitions and mathematical statement}
Fix a finite-signature reduct $A$ of a countable homogeneous relational structure $B$, where homogeneity extends every finite partial isomorphism to an automorphism, and finite boundedness describes the finite induced substructures by finitely many forbidden structures. The relations of $A$ are first-order definable in $B$. CSP$(A)$ takes a finite structure and asks for a homomorphism into the fixed template $A$. A primitive-positive definition uses existential quantifiers and conjunctions of atomic formulas. In the source's construction convention, primitive-positive powers and homomorphic equivalence determine whether $A$ can construct $K_3$. The assertion is
\[
 K_3\text{ has no such pp-construction from }A
 \quad\Longrightarrow\quad\operatorname{CSP}(A)\in\mathrm P.
\]
The template is fixed, not part of the finite input; its precise pp-construction convention is retained from the source.
\par\smallskip\textit{Formulation and concept references:} \href{https://arxiv.org/html/2502.06621v4}{[S1]}.

\subsection{Short English statement}
For the specified infinite-domain templates, does the absence of a primitive-positive construction of three-coloring guarantee a polynomial-time constraint-satisfaction algorithm?
\subsection{Sources}
\begin{itemize}
\item[S1]Pinsker, Rydval, Schöbi, Spiess and Winkler, Three Fundamental Questions in Infinite-Domain Constraint Satisfaction, arXiv:2502.06621v4 (6 May 2026).
\url{https://arxiv.org/html/2502.06621v4}.
\textit{Formulation source.}
\end{itemize}
\end{document}
